% ============================================
% Assistant Professor Cover Letter – University of Illinois Chicago
% ============================================

\documentclass[11pt,letterpaper]{letter}

\usepackage{geometry}
\usepackage{microtype}
\usepackage[hidelinks]{hyperref}

\geometry{left=25mm,right=25mm,top=25mm,bottom=25mm}

\signature{Decory Edwards}
\address{
Department of Economics \\
Johns Hopkins University \\
Baltimore, MD 21218 \\
\texttt{dedwar65@jh.edu}
}

\begin{document}

\begin{letter}{Search Committee \\
Department of Economics \\
University of Illinois Chicago \\
Chicago, IL \\ USA}

\opening{Dear Members of the Search Committee,}

I am writing to apply for the tenure-track Assistant Professor position in Economics at the University of Illinois Chicago. I am a Ph.D.\ candidate in Economics at Johns Hopkins University, advised by Professors Christopher D.~Carroll, Nicholas W.~Papageorge, and Stelios Fourakis, and I expect to receive my degree in May~2026. My research fields are macroeconomics, wealth and income dynamics, and computational economics.

My research agenda focuses on understanding the determinants of wealth inequality and the mechanisms through which heterogeneity in returns to wealth shapes aggregate outcomes. In my job-market paper, I estimate the degree of heterogeneity in individual portfolio returns using micro-data from the Survey of Consumer Finances and simulate a structural model in which households face idiosyncratic return risk. I show that allowing for persistent heterogeneity in returns substantially improves the model’s ability to match the empirical distribution of wealth, offering a tractable way to connect micro-level return dispersion to macroeconomic inequality.

In a second project, I use the Health and Retirement Study to examine the relationship between interpersonal trust and portfolio performance. Building on the literature that finds a hump-shaped relationship between trust and income, I show that a similar relationship holds for individual returns to wealth measured in the HRS data, highlighting a behavioral channel through which social attitudes shape saving and investment outcomes. This work also provides new estimates of the persistent component of returns to net worth, which motivate the calibration of heterogeneity in my structural model.

The University of Illinois Chicago’s public economics focus is a particularly strong fit for my research interests. UIC’s commitment to rigorous applied research and its role within a diverse urban research university provide an ideal environment for work at the intersection of quantitative methods and policy relevance. I am especially excited by the opportunity to engage with faculty who are actively publishing in public and applied economics and to contribute to UIC’s growing presence in these areas.

\textbf{Teaching.} I have experience teaching and supporting courses in \textit{Principles of Macroeconomics}, \textit{Intermediate Macroeconomic Theory}, and applied quantitative economics. My teaching emphasizes economic intuition, empirical grounding, and connections between theory and policy. At UIC, I would be well prepared to teach both undergraduate and Ph.D. courses, including core theory, public economics, and quantitative methods. I value inclusive pedagogy and am committed to engaging students from diverse academic and personal backgrounds.

I am enthusiastic about the possibility of contributing to UIC’s academic mission in research, teaching, and service. Thank you very much for your time and consideration. I would welcome the opportunity to discuss my application further.

\closing{Sincerely,}

\end{letter}
\end{document}
